%% kapa: exam=1 round=32 year=2021 period=1 form=A subject=civil_law subject_ko=민법 seq=18 number=18
%% source: Q-Net 감정평가사 기출문제 (cst00309) artlSeq=5208397 file=extracted/1008/5208397/A형/2021년도 제32회 감정평가사 1차 1교시 A형.pdf
%% conversion: kordoc 4.24.51 -> tools/gen1.py (auto); review=proofread-sonnet; flags=none
\begin{question}{18}
법률행위의 조건과 기한에 관한 설명으로 옳은 것은? (다툼이 있으면 판례에 따름)
\begin{choices}
  \item 법정조건도 법률행위의 부관으로서 조건에 해당한다.
  \item 채무면제와 같은 단독행위에는 조건을 붙일 수 없다.
  \item 기한은 특별한 사정이 없는 한 채권자의 이익을 위한 것으로 추정한다.
  \item 조건에 친하지 않은 법률행위에 불법조건을 붙이면 조건 없는 법률행위로 전환된다.
  \item 불확정한 사실의 발생을 기한으로 한 경우, 특별한 사정이 없는 한 그 사실의 발생이 불가능한 것으로 확정된 때에도 기한이 도래한 것으로 본다.
\end{choices}
\end{question}
\answer{5}
